% Source-derived chapter 04: Smoothness
% Original source: refined-brill-noether-theory-hurwitz-spaces
\chapter{Smoothness}
\label{refined-brill-noether-theory-hurwitz-spaces-chapter-04}
\source{refined-brill-noether-theory-hurwitz-spaces}

\begin{remark}[Source section]
\label{refined-brill-noether-theory-hurwitz-spaces-chapter-04-source}
\source{refined-brill-noether-theory-hurwitz-spaces}
\source{refined-brill-noether-theory-hurwitz-spaces}
This chapter reproduces the corresponding section of the retrieved
LaTeX source, including its definitions, statements, examples, and
proofs. It has not yet been translated into Lean.
\notready
\end{remark}

% The source section title was: Smoothness
\label{smoothness}
\source{refined-brill-noether-theory-hurwitz-spaces}

In this section, we prove that $\Sigma_{\vec{e}}(C, f)$ is smooth for general $f\colon  C \rightarrow \pp^1$. This should be thought of as an analogue of the Gieseker-Petri theorem, which was first proved by Gieseker \cite{G}, and later by Eisenbud-Harris \cite{EH2} using a degeneration with elliptic tails.

For every $L \in \Pic^d(C)$, there is a natural map
\begin{equation} \label{dmap}
\source{refined-brill-noether-theory-hurwitz-spaces}
H^1(\pp^1, f_*\O_C) =H^1(C, \O_C) = \Def^1(L) \rightarrow \Def^1(f_*L) = H^1(\pp^1, End(f_*L)),
\end{equation}
sending a first order deformation of $L$ to the induced deformation of the push forward. This map is realized by taking cohomology of the map of sheaves $\eta: f_*\O_C \rightarrow End(f_*L)$ that locally sends a function $z$ on $C$ to the endomorphism ``multiplication by $z$" on $L$, viewed as an $\O_{\pp^1}$ module.
The kernel of \eqref{dmap} is the tangent space to $\Sigma_{\vec{e}}(C, f)$.
Thus, our goal is to show that \eqref{dmap} is surjective for all $L \in \Pic^d(C)$. Indeed, this implies that if $\Sigma_{\vec{e}}(C, f)$ is non-empty, 
\[\dim T_L \Sigma_{\e}(C, f) \leq g - \dim \Def(f_*L) = g - u(\vec{e}) \leq \dim \Sigma_{\vec{e}}(C, f),\]
so it is smooth.

We proceed by showing that the Serre dual of \eqref{dmap},
\begin{equation} \label{sd}
\source{refined-brill-noether-theory-hurwitz-spaces}
\mu: H^0(\pp^1, End(f_*L) \otimes \omega_{\pp^1}) \rightarrow H^0(\pp^1, (f_*\O_C)^\vee \otimes \omega_{\pp^1}),
\end{equation}
is injective. The kernel of this map is the ``obstruction to smoothness." We think of $H^0(\pp^1, End(f_*L) \otimes \omega_{\pp^1})$ as the subspace of $H^0(\pp^1, End(f_*L))$ vanishing at two prescribed points. The map $\mu$ is thus a restriction of the map on global sections induced by
\[\tilde{\mu}: End(f_*L) \cong End(f_*L)^\vee \to (f_* \O_C)^\vee,\]
which is the composition of the canonical isomorphism $End(f_*L) \cong End(f^*L)^\vee$ with the map dual to $\eta$. For any vector bundle $E$, this isomorphism
$End(E) \cong End(E)^\vee$
 is induced by the perfect pairing $End(E) \times End(E) \to \O$ given by $(\phi, \psi) \mapsto \tr(\phi \cdot \psi)$. Therefore, $\tilde{\mu}$ sends
an endomorphism $\phi \in End(f_*L)(U)$ to the linear functional on $(f_*\O_C)(U)$ given by $z \mapsto \tr(\phi  \cdot \eta(z))$. We will need to know that this map is non-zero on certain elements over components of our degeneration.

\begin{Lemma}
\source{refined-brill-noether-theory-hurwitz-spaces}
\notready \label{nz}
\source{refined-brill-noether-theory-hurwitz-spaces}
Let $f\colon  X \rightarrow \pp^1$ be a degree $k$ map of an elliptic curve to $\pp^1$ and let $L \in \Pic^d(C)$.
If $\phi \in H^0(\pp^1, End(f_*L))$ is represented by a matrix with a single nonzero entry, then $\tilde{\mu}(\phi) \neq 0$. 
\end{Lemma}
\begin{proof}
For each open subset $U \subset \pp^1$, we have a commutative diagram
\begin{center}
\begin{tikzcd}
& H^0(\pp^1, End(f_*L)) \arrow{d} \arrow{r}{\tilde{\mu}} & \arrow{d} H^0(\pp^1, (f_*\O_C)^\vee) \\
& H^0(U, End(f_*L)) \arrow{r}[swap]{(\tilde{\mu})|_U} & H^0(U, (f_*\O_C)^\vee).
\end{tikzcd}
\end{center}
It suffices to show that the image of $\phi$ in the lower right is nonzero.
Choose $U$ small enough that $L$ is trivialized on $f^{-1}(U)$ and $f_*\O_C$ is trivialized on $U$, so
we have isomorphisms $(f_*L)|_U \cong (f_*\O_C)|_U \cong \O_{U}^{\oplus k}$. By hypothesis, there exists a basis so that, $\phi|_U : \O_U^{\oplus k} \rightarrow \O_U^{\oplus k}$ is represented by a matrix with one non-zero entry, say in the $(i, j)$ slot. These basis vectors of $\O_U^{\oplus k}$ correspond to non-vanishing functions in $\O_C(f^{-1}(U))$. The ratio of the $j$th basis element over the $i$th basis element therefore defines a function $z \in \O_C(f^{-1}(U))$ such that the $(j, i)$ entry of $\eta(z)$ is non-zero.
Hence, $\tr(\phi|_U \cdot \eta(z)) \neq 0$, showing $\tilde{\mu}(\phi|_U)$ is non-zero.
\end{proof}

We now deduce the desired injectivity by studying limits on the central fiber of our degeneration from the previous section, continuing all notation developed there.
Recall that for each $\L_0 \in \Pic^d(\X_0)$, we understood $V_\L \subset H^0(End(\f_*\L_0))$ through its image under the inclusion $\iota$ as compatible tuples $(\phi_1, \ldots, \phi_g)$ in $\bigoplus_{i=1}^g H^0(P_i, End(E_i))$.
\begin{Lemma}
\source{refined-brill-noether-theory-hurwitz-spaces}
\notready \label{sm}
\source{refined-brill-noether-theory-hurwitz-spaces}
For general $\X_t \rightarrow \P_t $ in our degeneration,
\[
\mu_t: H^0(\P_t, End(f_*\L_t) \otimes \omega_{\P_t}) \rightarrow H^0(\P_t, (f_*\O_{\X_t})^\vee \otimes \omega_{\P_t}),
\]
is injective for all $\L_t \in \Pic^d(\X_t)$. Hence, if it is non-empty, $\Sigma_{\vec{e}}(C, f)$ is smooth for general degree $k$ covers $f\colon  C \rightarrow \pp^1$.
\end{Lemma}
\begin{proof}
Let $\omega$ be the relative dualizing sheaf of $\P \rightarrow \Delta$. 
 Recall that $p_0$ and $p_g$ are points of total ramification of $\f|_{\X_0}$ that are distinct from the nodes. We set $\zeta_1 = f(p_0) \in P_1$ and $\zeta_g = f(p_g) \in P_g$, so we have an isomorphism $\omega|_{\P_0} \cong \O_{\P_0}(-\zeta_1 - \zeta_g)$.

Let $\L$ be given and define $V_\L(-\zeta_1-\zeta_g) \subset V_\L$ to be the subspace of sections vanishing at $\zeta_1$ and $\zeta_g$.
We have a commutative diagram
\begin{center}
\begin{tikzcd}
V_\L(-\zeta_1-\zeta_g) \arrow[swap]{ddr}{\iota} \arrow{r} &H^0(\P_0, End(\f_*\L_0) \otimes \omega|_{\P_0}) \arrow{d} \arrow{r}{\mu} &H^0(\P_0, (\f_*\O_{\X_0})^\vee \otimes \omega|_{\P_0}) \arrow{d} \\
&H^0(\P_0, End(\f_*\L_0)) \arrow{d} \arrow{r}{\tilde{\mu}} &H^0(\P_0, (\f_*\O_{\X_0})^\vee) \arrow{d}\\
 &\bigoplus_{i=1}^g H^0(P_i, End(E_i)) \arrow[swap]{r}{\oplus \tilde{\mu}_i} &\bigoplus_{i=1}^g H^0(P_i, ((f_i)_*\O_{X_i})^\vee).
\end{tikzcd}
\end{center}
By uppersemi-continuity, injectivity of $\mu_t$ for general $t$ follows from showing the composition along the top row is injective. We will show that the composition from the upper left to the lower right along the bottom is injective.

For each $i$, 
let $W_p^i \subset H^0(P_i, End(E_i))$ denote the subspace of endomorphisms on component $i$ that are order preserving at $p$. In addition, let $\diag^i_{p} : W_{p}^i \rightarrow \cc^{\oplus k}$ be defined by $\phi_i \mapsto (d_{p}^{(0)}(\phi_i), \ldots, d_{p}^{(k-1)}(\phi_i))$, which we saw in Lemma \ref{zing} corresponds to taking diagonal entries of a matrix representative for an endomorphism. 
On $W_{p_{i-1}} \cap W_{p_{i}}$, the maps $\diag^i_{p_i}$ and $\diag^i_{p_{i-1}}$ are both defined  and are related by a permutation (they correspond to taking diagonal entries of a matrix in different orders). Hence,
\begin{equation} \label{twoeq}
\source{refined-brill-noether-theory-hurwitz-spaces}
W_{p_{i-1}} \cap W_{p_{i}} \cap \ker(\diag^i_{p_i}) = W_{p_{i-1}} \cap W_{p_{i}} \cap \ker(\diag^i_{p_{i-1}}).
\end{equation}
Lemma \ref{cond} \eqref{d}, implies that in a compatible tuple, if $\phi_{i-1} \in \ker(\diag^{i-1}_{p_i})$ then $\phi_i \in \ker(\diag^i_{p_i})$.
Note that, taking a matrix representation for $\phi_1$ as in \eqref{mat}, the condition $\phi_1(\zeta_1) = 0$ is that $c_{j, \ell} = 0$ and $\alpha_{\ell, j}  = 0$ for \textit{all} $\ell, j$. Thus, if $\phi_1(\zeta_1) = 0$, then $\phi_1 \in  W_{p_0}^1 \cap \ker(\diag^1_{p_0})$. 
Similarly, if $\phi_g(\zeta_g) = 0$ then $\phi_g \in W_{p_g} \cap \ker(\diag^g_{p_g})$.
By Lemma \ref{cond} and \eqref{twoeq}, we therefore have
\begin{equation*}
\iota(V_\L(-\zeta_1-\zeta_g)) \subseteq \left\{(\phi_1, \ldots, \phi_g) : 
\text{ $\phi_i \in W_{p_{i-1}}^i \cap W_{p_i}^i \cap \ker(\diag^{i}_{p_{i}})$ for $1 \leq i \leq g$}
\right\}.
\end{equation*}
If $g > 1$, then we can choose a degree distribution so that $\deg(L_i)$ is never a multiple of $k$, and hence $L_i$ is never $f^*\O(n)$.
The final sentence of Lemma \ref{zing} then ensures that each $\phi_i$ is represented by a matrix with at most one-nonzero entry. In the case $g = 1$ we need an additional argument if $L_1 = f^*\O(n)$. In this case, choosing any splitting of $f_*L_1$ induces a splitting of $End(f_*L_1)$ where all but one of the matrix entries consist of a constant or linear form, and one entry is quadratic. After imposing vanishing at $\zeta_1$ and $\zeta_g$, only the quadratic entry can be non-zero. In either case, all $\phi_i$ have at most one non-zero entry, so Lemma \ref{nz} now shows that the composition of the inclusion $\iota$ with $\oplus \tilde{\mu}_i$ is injective.
\end{proof}
